\section{因子分析モデルの行列表現}
\subsection{授業内容}
\subsubsection{科目の中でのこのコマの位置づけ}
因子負荷量を算出するためには，線形代数でいうところの固有値についての理解が必要である。
まずは固有値と固有ベクトルを導入し，どのように計算するかを見る。特に固有ベクトルはノルムが定まらないことを確認する。
そこから，固有値と固有ベクトルがどのような性質を持っているかを幾何学的観点から確認する。正方行列が座標変換を行うためのものであると考えれば，固有ベクトルは変換行列の基底となることがわかるだろう。データ解析にあたって，相関行列の基底を求めるとはどういうことかをイメージするだけでも，因子分析の理解がまた一歩深まるだろう。
\subsubsection{コマ主題細目}
\begin{description}
    \item[固有値と固有ベクトル] 行列の固有値と固有ベクトルの性質を理解する。特にトレースが対角要素の和になることは，線形代数的観点からは重要な気づきを得る。
    \item[固有ベクトルを求める] $2\times2$行列を例に，固有値と固有ベクトルを求める計算を行う。固有方程式を解くことになるが，$n$次元に一般化するのは難しく，矜持的に計算することを知る。
    \item[固有値と固有ベクトルの幾何学的意味] 正方行列は一次変換行列であり，固有ベクトルはその基底であることを単純な行列から理解する。
    \item[因子分析モデルの意味] 因子分析モデルは相関行列を固有値分解することであり，それはすなわち相関行列の中にある基本的な次元・座標を求めることにある。すなわち複数人の反応パターンの共通要素を取り出すということであり，これは心理学的アプローチをほぼ直接的に数学表現したものであることを理解する。座標の回転についても触れ，仮定を緩めた場合の表現も理解する。
\end{description}
\subsubsection{キーワード}
\begin{itemize}
\item 因子分析モデルの行列表現
\item 固有値
\item 固有ベクトル
\item 固有ベクトルの幾何学的理解
\end{itemize}
\subsection{授業情報}
\paragraph{コマの展開方法}
講義
\subsubsection{予習・復習課題}
\paragraph{予習}因子分析の基本モデル，第一・第二定理の導出を復習しておこう。
\paragraph{復習}因子分析モデルが何をやっているかを考えた上で，心理学における尺度の利用やその解釈においてどのような注意をしなければならないかを言語化してみよう。
